\documentclass[12pt, letterpaper]{article}

\usepackage[utf8]{inputenc}
\usepackage[T1]{fontenc}
\usepackage[spanish]{babel}
\usepackage{geometry}
\usepackage{amsmath, amssymb, amsthm}
\usepackage{booktabs}
\usepackage{array}
\usepackage{microtype}
\usepackage{setspace}
\usepackage{parskip}
\usepackage{titlesec}
\usepackage{fancyhdr}
\usepackage{enumitem}
\usepackage{bm}
\usepackage{mathtools}
\usepackage{tcolorbox}
\usepackage{hyperref}

\tcbuselibrary{skins, breakable}

\geometry{top=2.8cm, bottom=3.0cm, left=3.2cm, right=3.2cm, headheight=14pt}
\setstretch{1.20}

\pagestyle{fancy}
\fancyhf{}
\fancyhead[C]{\small\textsc{Econometría --- Gujarati --- Ejercicio 13.19}}
\fancyfoot[C]{\small\thepage}
\renewcommand{\headrulewidth}{0.4pt}
\renewcommand{\footrulewidth}{0pt}

\titleformat{\section}
  {\large\bfseries\scshape}{\thesection.}{0.6em}{}
  [\vspace{-4pt}\rule{\linewidth}{0.4pt}]
\titleformat{\subsection}{\normalsize\bfseries}{\thesubsection.}{0.5em}{}
\titleformat{\subsubsection}{\normalsize\itshape}{}{0em}{}

\newtcolorbox{clave}{
  colback=white, colframe=black,
  boxrule=0.5pt, arc=3pt,
  left=8pt, right=8pt, top=6pt, bottom=6pt,
  breakable
}

\newcommand{\E}{\operatorname{E}}
\newcommand{\Cov}{\operatorname{Cov}}
\newcommand{\Var}{\operatorname{Var}}

\begin{document}

\begin{center}
  {\LARGE\bfseries La Prueba RESET de Ramsey}\\[6pt]
  {\LARGE\bfseries y su Equivalencia con el Modelo Cuadrático Directo}\\[12pt]
  {\large\itshape Demostración de que detectar $\hat{\alpha}_3 \neq 0$ en RESET}\\
  {\large\itshape equivale a estimar $Y_i = \beta_1 + \beta_2 X_i + \beta_3 X_i^2 + u_i$}\\[14pt]
  \rule{0.55\linewidth}{0.6pt}\\[8pt]
  {\normalsize Ejercicio~13.19 --- \textit{Econometría}, Gujarati \& Porter, 5.a ed.}\\[4pt]
  {\small\textsc{Demostración conceptual, teórica y algebraica}}
\end{center}

\vspace{10pt}

%------------------------------------------------------------------------------
\section{Marco conceptual y teórico}
%------------------------------------------------------------------------------

\subsection{La prueba RESET: motivación}

La prueba RESET (\textit{Regression Equation Specification Error Test}),
propuesta por Ramsey (1969), es una herramienta para detectar
\textbf{errores de especificación funcional}: en particular, la omisión
de términos no lineales (potencias o productos cruzados de los regresores)
en un modelo de regresión.

La idea central es la siguiente. Si el modelo lineal
\begin{equation}
  Y_i = \beta_1 + \beta_2 X_i + u_i
  \label{eq:original}
\end{equation}
está \emph{correctamente especificado}, entonces ninguna transformación
no lineal de los regresores debería tener poder explicativo adicional
sobre $Y_i$, una vez que $X_i$ ha sido incluida. En particular, las
potencias de los valores ajustados $\hat{Y}_i^2$, $\hat{Y}_i^3$, etc.,
no deberían ser significativas.

\subsection{¿Por qué usar potencias de $\hat{Y}_i$ en lugar de $X_i^2$ directamente?}

Dado que $\hat{Y}_i = \hat{\beta}_1 + \hat{\beta}_2 X_i$ es una
combinación lineal de los regresores originales, sus potencias son
combinaciones de términos polinomiales en $X_i$:
\[
  \hat{Y}_i^2 = (\hat{\beta}_1 + \hat{\beta}_2 X_i)^2
  = \hat{\beta}_1^2 + 2\hat{\beta}_1\hat{\beta}_2 X_i
    + \hat{\beta}_2^2 X_i^2.
\]
Así, $\hat{Y}_i^2$ contiene información sobre $X_i^2$; $\hat{Y}_i^3$
contiene información sobre $X_i^3$; y en general, $\hat{Y}_i^p$ captura
términos polinomiales hasta el orden $p$. La ventaja operativa es que
se resume en \emph{una sola} variable adicional la información que en
principio requeriría añadir múltiples potencias de $X_i$, preservando
así los grados de libertad.

\subsection{Hipótesis de la prueba}

La prueba RESET contrasta:
\[
  H_0: \alpha_3 = 0
  \qquad \text{(forma funcional correcta)}
  \qquad \text{vs.} \qquad
  H_1: \alpha_3 \neq 0
  \qquad \text{(mala especificación funcional)}.
\]
Si $H_0$ se rechaza, el modelo lineal~\eqref{eq:original} es inadecuado
y se debe explorar una forma funcional más rica, en particular la
inclusión del término $X_i^2$.

%------------------------------------------------------------------------------
\section{Planteamiento formal del procedimiento RESET}
%------------------------------------------------------------------------------

El procedimiento se desarrolla en tres etapas:

\begin{enumerate}[leftmargin=1.8cm, label=\textbf{Etapa \arabic*.}]

  \item \textbf{Estimar el modelo original por MCO.} Se obtienen los
    estimadores $\hat{\beta}_1$ y $\hat{\beta}_2$, y los valores ajustados:
    \begin{equation}
      \hat{Y}_i = \hat{\beta}_1 + \hat{\beta}_2 X_i.
      \label{eq:yhat}
    \end{equation}

  \item \textbf{Estimar la regresión aumentada (RESET).} Se agrega el
    cuadrado de los valores ajustados como regresor adicional:
    \begin{equation}
      Y_i = \alpha_1 + \alpha_2 X_i + \alpha_3 \hat{Y}_i^2 + v_i.
      \label{eq:reset}
    \end{equation}

  \item \textbf{Contrastar la significancia de $\hat{\alpha}_3$.}
    Se usa la razón $t$ o la prueba $F$ para evaluar si $\alpha_3 = 0$.
    Un rechazo indica mala especificación funcional en el modelo original.

\end{enumerate}

%------------------------------------------------------------------------------
\section{Demostración algebraica de la equivalencia}
%------------------------------------------------------------------------------

\subsection{Paso 1 --- Expresión explícita de $\hat{Y}_i^2$}

Partiendo de la definición~\eqref{eq:yhat}, elevamos al cuadrado:
\begin{align}
  \hat{Y}_i^2
  &= \bigl(\hat{\beta}_1 + \hat{\beta}_2 X_i\bigr)^2 \notag\\
  &= \hat{\beta}_1^2
     + 2\hat{\beta}_1\hat{\beta}_2\, X_i
     + \hat{\beta}_2^2\, X_i^2.
  \label{eq:yhat2}
\end{align}
Esta expansión es clave: muestra que $\hat{Y}_i^2$ es una
\emph{función exacta} de $1$, $X_i$ y $X_i^2$, con coeficientes que
dependen exclusivamente de los estimadores MCO del modelo original.

\subsection{Paso 2 --- Sustitución en la regresión RESET}

Insertamos~\eqref{eq:yhat2} en el modelo aumentado~\eqref{eq:reset}:
\begin{align}
  Y_i
  &= \alpha_1 + \alpha_2 X_i
     + \alpha_3\Bigl[\hat{\beta}_1^2
       + 2\hat{\beta}_1\hat{\beta}_2 X_i
       + \hat{\beta}_2^2 X_i^2\Bigr]
     + v_i.
  \label{eq:reset_expandido}
\end{align}

\subsection{Paso 3 --- Distribución y agrupamiento por potencias de $X_i$}

Distribuimos $\alpha_3$ y agrupamos los términos según su potencia
en $X_i$:
\begin{align}
  Y_i
  &= \underbrace{\bigl(\alpha_1 + \alpha_3\hat{\beta}_1^2\bigr)}_{\displaystyle\beta_1^*}
   + \underbrace{\bigl(\alpha_2 + 2\alpha_3\hat{\beta}_1\hat{\beta}_2\bigr)}_{\displaystyle\beta_2^*}
     X_i
   + \underbrace{\alpha_3\hat{\beta}_2^2}_{\displaystyle\beta_3^*}
     X_i^2
   + v_i.
  \label{eq:agrupado}
\end{align}

\subsection{Paso 4 --- Identificación de los nuevos parámetros}

Definimos los coeficientes reparametrizados como:
\begin{align}
  \beta_1^* &= \alpha_1 + \alpha_3\hat{\beta}_1^2,
  \label{eq:beta1star}\\
  \beta_2^* &= \alpha_2 + 2\alpha_3\hat{\beta}_1\hat{\beta}_2,
  \label{eq:beta2star}\\
  \beta_3^* &= \alpha_3\hat{\beta}_2^2.
  \label{eq:beta3star}
\end{align}
Con estas definiciones, la ecuación~\eqref{eq:agrupado} se escribe como:
\begin{equation}
  \boxed{Y_i = \beta_1^* + \beta_2^* X_i + \beta_3^* X_i^2 + v_i,}
  \label{eq:cuadratico}
\end{equation}
que es precisamente el modelo cuadrático directo que el ejercicio pide
demostrar. \hfill $\blacksquare$

%------------------------------------------------------------------------------
\section{Interpretación de la equivalencia}
%------------------------------------------------------------------------------

\subsection{Equivalencia entre $\hat{\alpha}_3 \neq 0$ y $\hat{\beta}_3^* \neq 0$}

La relación~\eqref{eq:beta3star} establece:
\[
  \beta_3^* = \alpha_3 \cdot \hat{\beta}_2^2.
\]
Dado que $\hat{\beta}_2^2 > 0$ siempre que $\hat{\beta}_2 \neq 0$
(es decir, siempre que $X_i$ sea relevante en el modelo original),
se cumple la equivalencia:
\begin{equation}
  \alpha_3 \neq 0
  \quad \Longleftrightarrow \quad
  \beta_3^* \neq 0.
  \label{eq:equivalencia}
\end{equation}
Por tanto, \textbf{rechazar $H_0: \alpha_3 = 0$ en la regresión RESET
es exactamente equivalente a rechazar $H_0: \beta_3^* = 0$ en el modelo
cuadrático directo}~\eqref{eq:cuadratico}. La prueba RESET detecta la
omisión de $X_i^2$ de manera \emph{indirecta} a través del cuadrado
de los valores ajustados, sin necesidad de especificar \textit{a priori}
qué término no lineal fue omitido.

\subsection{Lo que la significancia de $\hat{\alpha}_3$ revela}

Si $\hat{\alpha}_3$ es estadísticamente significativa en la regresión
aumentada~\eqref{eq:reset}, se concluye que:

\begin{enumerate}[leftmargin=1.6cm, label=\textbf{\arabic*.}]
  \item El modelo lineal~\eqref{eq:original} está \textbf{mal especificado}:
    omite al menos el término cuadrático $X_i^2$.
  \item La relación verdadera entre $Y_i$ y $X_i$ es \textbf{no lineal},
    y puede ser aproximada por el modelo cuadrático~\eqref{eq:cuadratico}.
  \item Los estimadores MCO del modelo original ($\hat{\beta}_1$,
    $\hat{\beta}_2$) son \textbf{sesgados e inconsistentes} debido a la
    omisión del término $X_i^2$.
\end{enumerate}

\subsection{Lo que la no significancia de $\hat{\alpha}_3$ implica}

Si $\hat{\alpha}_3$ no resulta significativa, no hay evidencia en los
datos de que el modelo lineal sea inadecuado \emph{frente a la alternativa
cuadrática}. Sin embargo, es importante notar que la prueba RESET con
$\hat{Y}_i^2$ únicamente contrasta contra alternativas polinomiales
de segundo grado; no detecta todas las posibles formas de mala
especificación (por ejemplo, relaciones logarítmicas o exponenciales
que no sean bien aproximadas por un polinomio de segundo orden).

\subsection{Comparación entre el procedimiento RESET y la estimación directa}

\begin{center}
\renewcommand{\arraystretch}{1.4}
\small
\begin{tabular}{p{4.0cm} p{4.5cm} p{4.5cm}}
\toprule
\textbf{Criterio} & \textbf{Modelo RESET} & \textbf{Modelo cuadrático directo} \\
\midrule
Variables explicativas
  & $X_i$ y $\hat{Y}_i^2$
  & $X_i$ y $X_i^2$ \\
Hipótesis contrastada
  & $H_0: \alpha_3 = 0$
  & $H_0: \beta_3^* = 0$ \\
Equivalencia
  & \multicolumn{2}{c}{$\alpha_3 \neq 0 \;\Leftrightarrow\; \beta_3^* \neq 0$
    (siempre que $\hat{\beta}_2 \neq 0$)} \\
Ventaja operativa
  & Una sola variable adicional resume $X_i^2$ sin especificarla \textit{a priori}
  & Requiere conocer \textit{a priori} que $X_i^2$ es el término omitido \\
\bottomrule
\end{tabular}
\end{center}

\begin{clave}
  \textbf{Conclusión:} La sustitución algebraica de $\hat{Y}_i^2$ en la
  regresión RESET revela que el término adicional no es más que una
  combinación lineal de $1$, $X_i$ y $X_i^2$. Al agrupar los coeficientes
  se obtiene exactamente el modelo cuadrático
  $Y_i = \beta_1^* + \beta_2^* X_i + \beta_3^* X_i^2 + v_i$,
  con $\beta_3^* = \alpha_3\hat{\beta}_2^2$.
  Dado que $\hat{\beta}_2^2 > 0$, la significancia de $\alpha_3$ en el
  RESET es condición necesaria y suficiente para la significancia de
  $\beta_3^*$ en el modelo directo. La prueba RESET es, por tanto, una
  prueba \emph{implícita} de la relevancia de $X_i^2$, con la ventaja
  de no requerir que el investigador especifique de antemano qué término
  no lineal ha sido omitido.
\end{clave}

\end{document}
